%% =====================================================================
%%  B3-T-17-06 -- what B3 contributes to "The Slow Reform"
%%  -------------------------------------------------------------------
%%  LAYER A: APPEND-ONLY. Written once at the entry's write, then left
%%  alone. If the source has more to say later, add a bullet -- do not
%%  rewrite. Material found ONLY here goes in the `unique` slot with
%%  @unique set to yes, which makes it undeletable (rule R10).
%% =====================================================================
%% @kind: source
%% @id: B3-T-17-06
%% @book: B3
%% @topic: T-17-06
%% @locator: Law 45, pp. 373--382
%% @status: distilled
%% @unique: no
%% @integrated: yes
%% @updated: 2026-09-19

\dossier{B3}{T-17-06}{\bookshort{B3}}{Law 45, pp. 373--382}

\begin{distilled}
  \item Everyone understands the need for change in the abstract, but people are creatures of habit: too much innovation is traumatic and leads to revolt.
  \item If new to a position, make a show of respecting the old way of doing things.
  \item If change is necessary, make it feel like a gentle improvement on the past.
  \item Wolsey's fall: the minister who opposed the king's change was destroyed by it; the successful reformers appeal to precedent and sequence change below the revolt threshold.
\end{distilled}
